\section{Related Work}
\label{sec:related_work}
This section positions the proposed multi-agent framework for digital twins within four avenues of related work: digital twin definitions and scope boundaries, multi-agent system (MAS) engineering and design, conceptual relationships between agents and digital twins, and existing multi-agent digital twin systems.



\subsection{Multi-agent systems}
Multi-agent systems provide a design paradigm for complex software systems in which multiple autonomous entities coordinate and interact to achieve system-level objectives. \citeauthor{weyns2010architecture} presents an architecture-based approach to MAS engineering, arguing that software architecture and reusable architectural patterns are key for building MAS that remain maintainable and evolvable as complexity grows~\cite{weyns2010architecture}.
This line of work motivates treating agent roles, interaction structure, and the representation of the environment as first-class design concerns, rather than implementation details~\cite{weyns2010architecture}.

\citeauthor{kinny1997_modelling_design_mas} propose modeling techniques and a design methodology that adapt object-oriented modeling to the needs of agent systems, including explicit modeling of agent classes, responsibilities, services, interactions, and internal decision structure~\cite{kinny1997_modelling_design_mas}.
Their work focuses on enabling principled design and refinement of agent systems, particularly by supporting abstraction and decomposition of roles and responsibilities~\cite{kinny1997_modelling_design_mas}.
\citeauthor{zhao2004agent} provides a complementary architecture-oriented perspective in the context of interactive simulation systems, where agents are used to coordinate distributed components and manage interactions across modules~\cite{zhao2004agent}.
Across these sources, MAS research contributes design principles for modularity, separation of responsibilities, and coordination mechanisms, which are relevant when digital twins must integrate heterogeneous components while remaining adaptable over time~\cite{weyns2010architecture,kinny1997_modelling_design_mas,zhao2004agent}.

\subsection{UML-derived modeling languages for multi-agent systems}
While UML is widely adopted in software engineering, its standard metamodel offers limited support for key agent abstractions such as autonomy, social organization, and cognitive decision-making. \citeauthor{GuedesVicari2022UMLDerived} survey a range of UML-derived languages and profiles proposed within Agent-Oriented Software Engineering (AOSE) to better capture these MAS-specific concerns~\cite{GuedesVicari2022UMLDerived}. Their review spans early extensions such as AUML, which primarily enrich interaction modeling, as well as more comprehensive approaches that introduce explicit concepts for roles, organizations, environments, and mental state (e.g., AML, MAS-ML), plus domain-specific variants like SEA-ML for semantic-web agents and MASRML for requirements modeling~\cite{GuedesVicari2022UMLDerived}.

A recurring finding in the chapter is that support for cognitive agents, particularly BDI-style notions of beliefs, desires/goals, intentions, plans, perceptions, and their dynamics, remains uneven. Some notations capture social interaction and protocol structure well but provide weak or unclear mechanisms for representing mental state and its evolution, while more expressive approaches can lead to monolithic or diagram-heavy models and still leave ambiguity about how perceptions, belief updates, and plan selection are represented at the modeling level~\cite{GuedesVicari2022UMLDerived}. For this thesis, this body of work reinforces the need to make roles, protocols, and environment interaction explicit in the system design, while also motivating an architecture-level focus on clear responsibility boundaries and executable coordination mechanisms rather than relying on notation alone for end-to-end clarity~\cite{GuedesVicari2022UMLDerived}.


\subsection{Existing multi-agent digital twin systems}
Several systems demonstrate the feasibility of combining MAS with digital twins, particularly in manufacturing and logistics. \citeauthor{nie2023_cloud_edge_dt_mas} propose a multi-agent and cloud-edge orchestration framework for distributed production control. Their approach places agent-based control at both cloud and edge levels, uses industrial IoT data for real-time state updates, and employs the digital twin to reflect and optimize distributed production processes~\cite{nie2023_cloud_edge_dt_mas}.
This work illustrates how MAS can be used to distribute decision-making across organizational layers, while the digital twin provides the shared representation needed for coordination~\cite{nie2023_cloud_edge_dt_mas}.

\citeauthor{marah2024_madtwin} proposes a framework for multi-agent digital twin development and demonstrates it through a smart warehouse case study. The framework represents components of a cyber-physical system as autonomous physical agents paired with digital twin agents within a digital twin environment, with an emphasis on extensible modeling and integration~\cite{marah2024_madtwin}.
This line of work shows how agent abstractions can support autonomy and cooperation within a digital twin setting, while also highlighting the architectural effort involved in integrating physical and digital components in operational systems~\cite{marah2024_madtwin}.

% \subsection{Verification and validation for digital twins}
% As digital twins move from monitoring and visualization toward decision support and automated action, trustworthiness becomes a central requirement. Recent work emphasizes that credibility and V\&V concerns must be explicitly addressed, including the traceability of assumptions, the intended use, and evidence supporting predictive validity~\cite{shao2023credibility}. For digital twins in high-stakes settings, survey and perspective work highlights the importance of verification, validation, and uncertainty quantification (VVUQ) as a structured approach to ensuring safety and efficacy~\cite{sel2025_vvuq_precision_dt}. From a modeling and scaling perspective, work on predictive digital twins emphasizes probabilistic foundations and scalable updating, which also implies that uncertainty handling and model evidence become integral to the twin lifecycle~\cite{kapteyn2021probabilistic}.

% For multi-agent digital twin systems, V\&V raises an additional challenge: it is not sufficient to validate individual models in isolation if system-level behavior depends on agent coordination policies, dynamic task allocation, and emergent interactions. The literature motivates separating what is being validated (model fidelity, data assimilation, decision logic, coordination outcomes) and under what assumptions. This strengthens the motivation for explicit responsibility boundaries and well-defined interfaces, because they enable targeted evidence collection and clearer claims about what has been verified or validated for a specific purpose~\cite{boyes2022analysis,shao2023credibility}.


\subsection{Research gap}
\label{sec:research_gap}

The background and related work show broad agreement on the core building blocks of digital twin  systems, including continuous data ingestion from physical assets, a virtual model, analytics or simulation, and a feedback path to operations~\cite{shao2020framework,jones2020characterising,soori2023digital,inamdar2024_dt_microelectronics}. At the same time, the literature converges on the observation that current DT implementations struggle when they are scaled beyond narrowly scoped, single-system deployments, especially when systems must be extended incrementally, coupled across assets, and maintained over long lifecycles~\cite{Gil2024,Visser2024,Kapteyn2020,vrabic2018digital}.

Based on the provided sections, the research gap can be stated as a set of architecture-level shortcomings that remain insufficiently addressed when DTs are engineered for modularity, adaptability, and trustworthiness:

\paragraph{Insufficient architectural support for modular coupling and evolution}
Several works highlight limited platform support for coupling and ensuring consistency between coupled components in monolithic DT~\cite{Gil2024}. Architectural proposals for scalability and integration indicate that modularization is feasible and beneficial~\cite{Visser2024,Kapteyn2020}, but the background highlights that DTs still face recurring integration challenges across interoperability, synchronization, lifecycle maintainability, and governance when deployed in dynamic industrial environments~\cite{jones2020characterising,kritzinger2018digital,schluse2017experimentable,shao2019cybersecurity}.

\paragraph{A mismatch between existing integration mechanisms and the need for runtime adaptivity}
Component coupling standards and co-simulation tooling (for example, FMI-based model exchange and coupling libraries) support reusable composition and reliable time-synchronized execution, but they primarily assume predefined interaction structure, which limits adaptivity for evolving DT configurations~\cite{blochwitz2012fmi20,muscle3docs}. Service-oriented and workflow-based DT platforms improve deployment modularity and interoperability through independent services coordinated by explicit workflows~\cite{borodulin2017dtcloud,bellavista2024microservices_serverless_dt}. However, a workflow-centric integration style places system-level behavior in predefined orchestration logic, which makes runtime reorganization and goal-driven adaptation a distinct architectural problem rather than a built-in capability.


\paragraph{Limited consolidation on responsibility boundaries between DT artifacts and agent-based decision-making}
Conceptual work argues for a clear separation of concerns, where the DT is tightly tied to asset-facing state synchronization and interfaces, while agents operate in a broader context to provide distributed reasoning, optimization, and cross-asset coordination~\cite{mariani2022_dt_agents_cross_fertilisation}. System proposals demonstrate feasibility of multi-agent DTs in manufacturing settings such as distributed production control and smart warehouses~\cite{nie2023_cloud_edge_dt_mas,marah2024_madtwin}, and surveys summarize the potential of combining DTs and multi-agent systems (MAS) for adaptability and decentralized control~\cite{Pretel2024SLR,Kalyani2024Survey}. However, there is still a need for an architecture-level framework that makes such responsibility boundaries explicit and operational, in a way that supports modular integration of heterogeneous components while maintaining clear interfaces for coordination and evolution.


\paragraph{Trustworthiness methods that scale from model validation to multi-agent, system-level behavior.}
As DTs are increasingly used for decision support and automated action, the literature emphasizes credibility, verification and validation (V\&V), and uncertainty quantification as central requirements \cite{shao2023credibility,sel2025_vvuq_precision_dt,kapteyn2021probabilistic}. The related work further motivates that multi-agent DTs introduce an additional layer of V\&V difficulty, because system-level outcomes may depend on decision logic, coordination policies, task allocation, and emergent interactions rather than model fidelity alone \cite{boyes2022analysis,shao2023credibility}. What remains insufficiently resolved is how to structure a DT architecture so that validation targets can be separated (model, data assimilation, decision logic, coordination outcomes) and evidence can be gathered and traced to explicit architectural responsibilities in a modular, composable system.


The section above highlights the need for a domain-agnostic, modular DT architecture that supports consistent coupling and incremental evolution of DT components~\cite{Gil2024,Kapteyn2020,Visser2024}, goes beyond static coupling and workflow orchestration by enabling goal-driven, distributed coordination~\cite{Pretel2024SLR,mariani2022_dt_agents_cross_fertilisation}, and provides explicit responsibility boundaries and interfaces that make system-level V\&V practicable in multi-agent DT settings~\cite{shao2023credibility,sel2025_vvuq_precision_dt,boyes2022analysis}. This gap motivates an agent-based DT framework in which asset-facing DT capabilities and cross-asset reasoning are architecturally separated and organized so that modular integration and targeted assurance activities can be addressed together.

The proposed architecture targets this gap by using organizational structure as a primary element while allowing runtime modification driven by the agents. This supports coordination that remains systematic through roles and protocols, while still enabling reorganization such as forming new groups to address emergent conditions.

An MAS-based digital twin architecture intended to be domain-agnostic and adaptive, combining autonomy and coordination with explicit support for modular integration of heterogeneous components is proposed. By incorporating a formal organizational structure and separating agent reasoning from simulation interfacing, the approach targets limitations in prior MAS-for-digital twin systems and aims to provide a more modular and extensible platform.
